\section{Representations}

Let $\g$ be a semisimple Lie algebra,
$\g = \at + \bigoplus_{\alpha \in R} \g_\alpha = \at + \fn^+ + \fn^-$. Let $V$
be a finite-dimensional representation of $\g$. (We don't quite need
finite-dimensional; we'll see exactly what we need a bit later.)

\begin{proposition}\leavevmode
\begin{enumerate}
  \item $V = \bigoplus_{\lambda \in \at^*} V_\lambda$, where
    $V_\lambda = \setc{v \in V}{t(v) = \lambda(t)v \ \forall t \in \at}$ (this
    is the weight space decomposition with respect to $\at$).
  \item If $V_\lambda \neq 0$, then $\lambda(h_\alpha) \in \Z$ for all
    $\alpha \in R$. (Here, $h_\alpha$ is the element of the $(\sl_2)_\alpha$,
    $h_\alpha = \nu^{-1}(\alpha\cv)$.)
\end{enumerate}
\end{proposition}

\begin{proof}
Since $V$ is a finite-dimensional representation of $\g$, it is a
finite-dimensional representation of $(\sl_2)_\alpha$. Therefore, $h_\alpha$
acts diagonalisably on $V$ with eigenvalues $\lambda(h_\alpha) \in \Z$. The
first assertion follows because the $h_\alpha$ commute and span $\at$.
\end{proof}

\begin{definition}
$Q = \Z R = \bigoplus_{\alpha_i \in \Pi} \Z\alpha_i$ is the \emph{lattice of
roots}.
$P = \setc{\gamma \in \Q R}{(\gamma, \alpha\cv) \in \Z \ \forall \alpha \in R}$
(the dual lattice) is the \emph{lattice of weights}. Note that the condition
can be checked by computing $(\gamma, \alpha_i\cv)$ for $\alpha_i \in \Pi$.
\end{definition}

Since $(\beta, \alpha\cv) \in \Z$ for $\beta, \alpha \in R$, we see that
$Q \subseteq P$. Also, if $V$ is a finite-dimensional representation of $\g$
and $V_\lambda \neq 0$, then $\lambda \in P$, as
$\lambda(h_\alpha) = (\lambda, \alpha\cv)$.

\begin{exercise}\leavevmode
\begin{enumerate}
  \item $|P/Q| < \infty$. In fact, $|P/Q| = \det A$, where $A$ is the Cartan
    matrix.
  \item $W(P) \subseteq P$, where $W$ is the Weyl group.
\end{enumerate}
\end{exercise}

\begin{example}
In $\sl_2$, $R = \{\pm\alpha\}$, $Q = \Z\alpha$, $(\alpha, \alpha) = 2$, so
$P = \Z\alpha/2$, and $|P/Q| = 2 = \det(2)$.
\end{example}

\begin{definition}
For $V$ a finite-dimensional representation of $\g$, define
\[
  \ch V = \sum_{\lambda \in P} \dim V_\lambda \cdot e^\lambda \in \Z[P].
\]
Here, $e^\lambda$ is a formal symbol ($e^\lambda e^\mu = e^{\lambda+\mu}$)
forming a basis for $\Z[P]$.
\end{definition}

``People are looking at me as if I'm speaking some language I don't speak, let
alone you don't speak.''

\begin{example}
For $\sl_2$, $P = \Z\alpha/2$. Write $z = e^{\alpha/2}$; then
$\ch L(n) = z^n + \dots + z^{-n} = \frac{z^{n+1} - z^{-(n+1)}}{z - z^{-1}}$.
\end{example}

\begin{example}
Let $\g = \sl_3$, and $V = \g$ the adjoint representation. Then
\[
  \ch V = 2 + z + w + zw + z^{-1} + w^{-1} + z^{-1}w^{-1},
\]
where $w = e^{\alpha_1}$, $z = e^{\alpha_2}$. See figure~\ref{fig:sl3-adjoint}.
\end{example}

\begin{figure}[!htbp]
\centering
\begin{tikzpicture}[scale=1.25,every node/.append style={font=\small}]
  \foreach \i in {-2,...,2} \foreach \j in {-2,...,2} {
    \pgfmathsetmacro{\x}{\i + \j*cos(120)}
    \pgfmathsetmacro{\y}{\j*sin(120)}
    \pgfmathsetmacro{\r}{sqrt(\x*\x+\y*\y)}
    \ifdim \r pt < 2.1pt \node[dot] at (\x,\y) {}; \fi
  }
  \node[above right=0pt] at (0,0) {$2$};
  \foreach \a in {0,60,...,300} \node[above right=0pt] at (\a:1) {$1$};
  \node[below=3pt] at (0:1) {$\alpha_1$};
  \node[below=3pt] at (120:1) {$\alpha_2$};
  \node[below=3pt] at (60:1) {$\alpha_1 + \alpha_2$};
\end{tikzpicture}
\caption{Root lattice points of $\sl_3$ near the origin, with the dimensions
of the weight spaces of the adjoint representation. The six roots carry $1$,
the origin carries $2$, and the picture is $W$-invariant ($W = S_3$ here).}
\label{fig:sl3-adjoint}
\end{figure}

\begin{proposition}
If $V$ is a finite-dimensional representation of $\g$, then
$\dim V_\lambda = \dim V_{w\lambda}$ for all $w \in W$, i.e.\ $\ch V$ is
$W$-invariant, $\ch V \in \Z[P]^W$. (In particular, for $\sl_n$ this means that
the $\ch V$ are symmetric functions.)
\end{proposition}

We will sketch three proofs, of which the third one is the easiest, but the
first two reveal what is actually going on.

\begin{sketch}{1}
If $G$ is an algebraic group with Lie algebra $\g$ then (by a theorem)
$W = N(T)/T$ (where $T$ is the maximal torus in $G$ whose Lie algebra is
$\at$, and $N(T)$ is its normaliser).

\begin{example}
If $G = SL_n$ with $T$ the diagonal matrices, then $N(T)$ consists of the
matrices with exactly one nonzero entry in each row and column, and
$N(T)/T = S_n$, the permutation matrices.
\end{example}

That is, there exists $\dot w \in N(T)$ such that $\dot wT = w$. If $G$ acts on
$V$ (as it always does when $G$ is simply connected), then
$\dot w(V_\lambda) = V_{w\lambda}$, since
\[
  t\dot w(v) = \dot w(\dot w^{-1}t\dot w)v = \dot w\bigl(\lambda(\dot w^{-1}t\dot w)v\bigr).
\]
That is, $W$ is almost a subgroup of $G$. However, it is definitely not
actually a subgroup. E.g., in $SL_2$, we take
$\dot s = \bigl(\begin{smallmatrix} 0 & -1 \\ 1 & 0 \end{smallmatrix}\bigr)$,
and then
$\dot s^2 = \bigl(\begin{smallmatrix} -1 & 0 \\ 0 & -1 \end{smallmatrix}\bigr) \in T$
is not equal to $1$. That is, $W$ does not embed into $G$ and does not act on
$V$. (There is a small $2$-group, $\approx (\Z/2)^l$, that intervenes. This
$2$-group is also involved in determining the coefficients $c_{\alpha\beta}$.)
\end{sketch}

\begin{sketch}{2}
We can mimic the above argument in $\g$:
\[
  \begin{pmatrix} 0 & -1 \\ 1 & 0 \end{pmatrix}
  = \begin{pmatrix} 1 & 0 \\ 1 & 1 \end{pmatrix}
    \begin{pmatrix} 1 & -1 \\ 0 & 1 \end{pmatrix}
    \begin{pmatrix} 1 & 0 \\ 1 & 1 \end{pmatrix}
  = \exp(f)\exp(-e)\exp(f).
\]
Therefore, for each root $\alpha$ we define
$\dot s_\alpha = \exp(f_\alpha)\exp(-e_\alpha)\exp(f_\alpha)$, where the
matrix exponential is defined in the usual power series way.

\begin{exercise}\leavevmode
\begin{enumerate}
  \item If $V$ is a finite-dimensional representation of $\g$, then
    $e_\alpha$, $f_\alpha$ act nilpotently on $V$, so
    $\dot s_\alpha\colon V \to V$ is a well-defined finite sum.
  \item $\dot s_\alpha^2 = \epsilon_\alpha\colon V_\lambda \to V_\lambda$,
    where $\epsilon_\alpha^2 = 1$ and $\epsilon_\alpha$ acts by a constant.
    Determine this constant (in terms of $\lambda$).
  \item $\dot s_\alpha(V_\lambda) = V_{s_\alpha\lambda}$.
  \qedhere
\end{enumerate}
\end{exercise}
\end{sketch}

\begin{remark}
We did not really need $V$ to be finite-dimensional; what we need instead is
for each $e_\alpha$ and $f_\alpha$ to act \emph{locally nilpotently}. We say
$x\colon V \to V$ acts \emph{locally nilpotently} if for all $v \in V$ we have
$x^nv = 0$ for some $n$.
\end{remark}

\begin{exercise}
This is equivalent to saying that as an $(\sl_2)_\alpha$-module, $V$ splits
into a (possibly infinite) direct sum of finite-dimensional $\sl_2$-modules.
\end{exercise}

\begin{definition}
Such a $V$ is called \emph{integrable} (because you can ``integrate'' the Lie
algebra action to a group action).
\end{definition}

In the rest of the course, you can replace all instances of
``finite-dimensional'' by ``integrable'', and then all theorems and proofs will
work for the Kac--Moody algebras as well.

\begin{sketch}{3}
If we're just working with $\sl_2$ anyway, we shouldn't need all the other
stuff.

\begin{exercise}
Consider $V$ as a representation of $(\sl_2)_\alpha \times \at$. Then $V$
splits up into a direct sum of strings, each of which is of the form
\[
  \lambda \qquad \lambda - \alpha \qquad \lambda - 2\alpha \qquad \cdots \qquad \lambda - m\alpha,
\]
where $m = \lambda(h_\alpha)$ for the top weight $\lambda$ of the string.
\end{exercise}

Such a string is obviously $s_\alpha$-invariant, since $s_\alpha$ will just
flip the string from left to right (as the reflection did in $\sl_2$).
\end{sketch}

\lecture{16}

\begin{definition}
For $\lambda, \mu \in \at^*$ write $\mu \leq \lambda$ to mean
$\lambda - \mu = \sum k_i\alpha_i$ with $k_i \in \N$. The points
$\setc{\mu \in P}{\mu \leq \lambda}$ are the lattice points in an obtuse cone
(figure~\ref{fig:cone}).
\end{definition}

\begin{figure}[!htbp]
\centering
\begin{tikzpicture}[scale=0.9,every node/.append style={font=\small}]
  \fill[gray!20] (0,0) -- (180:2.8) -- ($(180:2.8)+(-60:2.5)$) -- (-60:2.5) -- cycle;
  \draw (0,0) -- (180:2.8);
  \draw (0,0) -- (-60:2.5);
  \node[dot] at (0,0) {};
  \node[above] at (0,0) {$\lambda$};
  \node[above] at (180:2.2) {$\lambda - \R_{\geq 0}\alpha_1$};
  \node[right] at (-60:2.2) {$\lambda - \R_{\geq 0}\alpha_2$};
\end{tikzpicture}
\caption{The cone $\setc{\mu}{\mu \leq \lambda}$ for $\sl_3$, bounded by the
directions $-\alpha_1$ and $-\alpha_2$.}
\label{fig:cone}
\end{figure}

\begin{definition}
If $V$ is a representation of $\g$, we say that $\mu \in P$ is a \emph{highest
weight} if $V_\mu \neq 0$ and $V_\lambda \neq 0 \implies \lambda \leq \mu$.

We say $v \in V_\gamma$ is a \emph{singular vector} if $v \neq 0$ and
$e_\alpha v = 0$ for all $\alpha \in R^+$. (Note that
$e_\alpha v \in V_{\gamma+\alpha}$, so if $\gamma$ is a highest weight, then
all nonzero elements of $V_\gamma$ are singular vectors.)

We say that $\mu$ is an \emph{extremal weight} if $w\mu$ is a highest weight
for some $w \in W$.
\end{definition}

\begin{example}
In the adjoint representation of $\sl_3$, the highest weight is
$\alpha_1 + \alpha_2$, and the six roots on the outer hexagon are the extremal
weights.
\end{example}

\begin{theorem}
Let $\g$ be a semisimple Lie algebra over $\C$.
\begin{enumerate}
  \item ``Complete reducibility'': if $V$ is a finite-dimensional
    representation of $\g$, then $V$ is a direct sum of irreducible
    representations of $\g$.

    Set $P^+ = \setc{\lambda \in P}{(\lambda, \alpha\cv) \geq 0 \ \forall \alpha \in R^+}$
    (of course, it's enough to check this on $\alpha \in \Pi$). This is the
    cone of \emph{dominant weights}. The assertion of the next parts of the
    theorem is that $P^+$ parametrises the finite-dimensional irreducible
    representations of $\g$, with $\lambda \mapsto L(\lambda)$. More
    precisely,
  \item Let $V$ be a finite-dimensional irreducible representation of $\g$,
    and $v \in V_\lambda$ a singular vector. Then
    \begin{alist}
      \item $\dim V_\lambda = 1$;
      \item $V_\mu \neq 0 \implies \mu \leq \lambda$, so $\lambda$ is a
        highest weight (and $v$ a highest-weight vector);
      \item $\lambda(h_i) \in \N$ for all $i$, i.e.\ $\lambda \in P^+$.
    \end{alist}
    Moreover, if $W$ is another finite-dimensional irreducible representation
    of $\g$ with highest weight $\lambda$, and $w \in W_\lambda$, $w \neq 0$,
    then there exists a unique isomorphism $V \to W$ mapping $v \mapsto w$.
    (Of course, given they are irreducible representations, if there is an
    isomorphism, there is a one-dimensional family of them, so we would pin
    down the exact one by sending $v \mapsto w$. That is, uniqueness is not
    the nontrivial statement here.)
  \item Given $\lambda \in P^+$, there exists a finite-dimensional irreducible
    representation with highest weight $\lambda$, called $L(\lambda)$.
  \item There is a formula for $\ch L(\lambda)$ (the Weyl character formula),
    to be stated later.
\end{enumerate}
\end{theorem}

\begin{remark}
For Kac--Moody algebras, we need a slight modification, since it is no longer
obvious that we will have highest-weight vectors. We need to work in the
category of integrable highest-weight representations instead.
\end{remark}

\begin{corollary}
\[
  \ch L(\lambda) = e^\lambda + \sum_{\mu < \lambda} a_{\mu\lambda}e^\mu \in \Z[P],
\]
hence $\setc{\ch L(\lambda)}{\lambda \in P^+}$ are linearly independent (they
have different leading terms). On the other hand,
\[
  \ch L(\lambda) = m_\lambda + \sum_{\substack{\mu < \lambda\\ \mu \in P^+}} \tilde a_{\mu\lambda}m_\mu
    \in \Z[P]^W, \qquad \text{where } m_\mu = \sum_{\gamma \in W\mu} e^\gamma.
\]
Since the $\{m_\mu\}_{\mu \in P^+}$ clearly form a basis for $\Z[P]^W$, the
$\ch L(\lambda)$ form a basis of $\Z[P]^W$.
\end{corollary}

\begin{corollary}
If $V$ and $W$ are finite-dimensional representations of $\g$, then
$V \cong W \iff \ch V = \ch W$.
\end{corollary}

\begin{proof}
Complete reducibility, plus the fact that the characters of irreducibles form a
basis.
\end{proof}

\begin{definition}
Define $\omega_i \in P$ to be the dual basis to the simple (co)roots,
$(\omega_i, \alpha_j\cv) = \delta_{ij}$ for $i, j = 1, \dots, l$. These are
called the \emph{fundamental weights}.
\end{definition}

Then $P^+ = \bigoplus \Z_{\geq 0}\omega_i = \setc{\sum n_i\omega_i}{n_i \geq 0}$.
This is a fundamental domain for the action of $W$ on $P$. Each irreducible
representation will have a unique highest weight lying in this cone, which is
what parametrises the representation.

\begin{exercise}
Compute the $\omega_i$ for the simple Lie algebras $\sl_n$, $\so_{2n}$,
$\so_{2n+1}$, $\spl_{2n}$, etc., and draw the picture for $A_2$, $B_2$, and
$G_2$.
\end{exercise}

\begin{figure}[!htbp]
\centering
\begin{tikzpicture}[scale=1.45,every node/.append style={font=\footnotesize}]
  \fill[gray!20] (0,0) -- (90:2.35) -- (30:2.35) -- cycle;
  \draw (0,0) -- (90:2.35); \draw (0,0) -- (30:2.35);
  \foreach \i in {-2,...,2} \foreach \j in {-2,...,2} {
    \pgfmathsetmacro{\x}{\i + \j*cos(120)}
    \pgfmathsetmacro{\y}{\j*sin(120)}
    \pgfmathsetmacro{\r}{sqrt(\x*\x+\y*\y)}
    \ifdim \r pt < 2.1pt \node[dot] at (\x,\y) {}; \fi
  }
  \node[below=2pt] at (0:1) {$\alpha_1$};
  \node[below=2pt] at (120:1) {$\alpha_2$};
  \node[above=2pt] at (60:1) {$\alpha_1 + \alpha_2$};
  \node[above] at (90:2.35) {$\R_{\geq 0}\,\omega_2$};
  \node[right] at (30:2.35) {$\R_{\geq 0}\,\omega_1$};
\end{tikzpicture}
\caption{Cone of dominant weights for $\sl_3$, spanned by the fundamental
weights. The adjoint representation's highest weight is $\alpha_1 + \alpha_2$.}
\label{fig:dominant-sl3}
\end{figure}

\begin{exercise}
$\lambda \in P \implies \lambda = \sum \lambda(h_i)\omega_i$, where
$h_i = \nu^{-1}(\alpha_i\cv)$.
\end{exercise}

\begin{example}
For any $\g$, the trivial representation is $L(0)$.
\end{example}

For any simple $\g$, the adjoint representation $V = \g$ is irreducible. The
decomposition $\g = \at + \bigoplus_{\alpha \in R} \g_\alpha$ means that a
highest weight $\lambda$ is a root such that $\lambda + \alpha_i \notin R$ for
all $i$, i.e.\ $\lambda = \theta$, the highest root in $R^+$ (see
figure~\ref{fig:dominant-sl3}). (Therefore, the theorem of the highest weight
implies that $\theta$ is unique, as promised.)

\begin{example}
In $A_{n-1}$ ($n \geq 3$), the highest root is $\theta = e_1 - e_n$ (recall
$\alpha_i = e_i - e_{i+1}$ and $h_i$ is a diagonal matrix with $+1$ in the
$i$th position, $-1$ in the $(i+1)$th position, and $0$ elsewhere). We have
$\theta(h_1) = 1$, $\theta(h_2) = 0, \dots, \theta(h_{n-2}) = 0$,
$\theta(h_{n-1}) = 1$. We can label the Dynkin diagram by the $\theta(h_i)$:
\[
\begin{tikzpicture}[baseline=-3pt,x=1.1cm]
  \node[node] (p1) at (0,0) {}; \node[node] (p2) at (1,0) {};
  \node[node] (p3) at (3,0) {}; \node[node] (p4) at (4,0) {};
  \draw (p1)--(p2); \draw[dotted] (p2)--(p3); \draw (p3)--(p4);
  \node[above=2pt] at (p1) {$1$}; \node[above=2pt] at (p2) {$0$};
  \node[above=2pt] at (p3) {$0$}; \node[above=2pt] at (p4) {$1$};
\end{tikzpicture}
\]
The labels give the decomposition $\theta = \omega_1 + \omega_{n-1}$ of
$\theta$ into fundamental weights.
\end{example}

\begin{exercise}
Compute the $\theta(h_i)$ for all the simple Lie algebras.
\end{exercise}

\begin{example}
Other representations of $\sl_n$: we have the standard representation,
$\C^n$, with weight basis $v_1, \dots, v_n$ and weights $e_1, \dots, e_n$
(recall $\sum e_i = 0$). The highest weight is $e_1$ (as
$e_1 > e_2 > \dots > e_n$, since $e_1 - e_2, e_2 - e_3, \dots \in R^+$).
Therefore, $\C^n = L(\omega_1)$, and
\[
  \ch \C^n = e^{e_1} + \dots + e^{e_n} = z_1 + \dots + z_n,
\]
where $z_i = e^{e_i} \in \Z[P] = \Z[z_1^{\pm 1}, \dots, z_n^{\pm 1}]/(z_1\cdots z_n = 1)$.
\end{example}

Now, if $V$ and $W$ are representations of $\g$, then so is $V \otimes W$ (via
$x \in \g$ acting by $x \otimes 1 + 1 \otimes x$). In particular,
$V \otimes V$ is a representation. Note that
$\sigma\colon V \otimes V \to V \otimes V$, $a \otimes b \mapsto b \otimes a$,
commutes with the $\g$-action, so the eigenspaces of $\sigma$ are
$\g$-modules. That is, $S^2V$ and $\Lambda^2V$ are $\g$-modules. In general,
these need not be irreducible, but they are for $V = \C^n$ and $\g = \sl_n$, as
we check below.

Consider $\Lambda^s\C^n$ (for $s \leq n - 1$). This has weight basis
$v_{i_1} \wedge \dots \wedge v_{i_s}$ for $i_1 < \dots < i_s$, and the weight
of this is $e_{i_1} + \dots + e_{i_s}$, as is clear from the action. Therefore,
\[
  E_i(v_{i_1} \wedge \dots \wedge v_{i_s}) = 0 \text{ for all } i
  \iff v_{i_1} \wedge \dots \wedge v_{i_s} = v_1 \wedge \dots \wedge v_s,
\]
i.e.\ this is the only singular vector (up to scalars). Note that its weight
is $\omega = e_1 + \dots + e_s$, which is the $s$th fundamental weight, as
$(\omega, e_i - e_{i+1}) = \delta_{is}$. Therefore,
\[
  \Lambda^s\C^n = L(\omega_s) \text{ is a \emph{fundamental representation}.}
\]
In particular, $\Lambda^{n-1}\C^n \cong (\C^n)^* = L(\omega_{n-1})$.

Consider $S^m\C^n$. It has weight basis $v_{i_1}\cdots v_{i_m}$ with
$i_1 \leq \dots \leq i_m$ and weight $e_{i_1} + \dots + e_{i_m}$. We have
\[
  E_i(v_{i_1}\cdots v_{i_m}) = 0 \text{ for all } i
  \iff v_{i_1}\cdots v_{i_m} = v_1\cdots v_1 = v_1^m,
\]
i.e.\ $S^m\C^n$ is an irreducible representation isomorphic to
$L(m\omega_1)$.

\begin{exercise}
Check that the above is true. Compute $\ch \Lambda^s\C^n$ and
$\ch S^m\C^n$. Find a closed formula for the generating functions
\[
  \sum_{m \geq 0} (\ch S^m\C^n)q^m, \qquad \sum_{m \geq 0} (\ch \Lambda^m\C^n)q^m.
\]
\end{exercise}

\lecture{17}

\begin{exercise}\leavevmode
\begin{enumerate}
  \item If $V$ and $W$ are $\g$-modules, $V$ finite-dimensional, then
    $V^* \otimes W \cong \Hom(V, W)$ as $\g$-modules, via
    $v^* \otimes w \mapsto (v \mapsto v^*(v) \cdot w)$.

    \begin{remark}
    If $V$ is a $\g$-module, $V^*$ is a $\g$-module with action
    $(x \cdot \lambda)(v) = -\lambda(xv)$. The $-$ sign comes from
    differentiating the group action, $(g \cdot \lambda)(v) = \lambda(g^{-1}v)$.
    \end{remark}
  \item Show that if $V = \C^n$ and $\g = \sl_n$, then
    $V \otimes V^* \cong \sl_n \oplus \C$, adjoint plus trivial
    representations (where the trivial representation comes from
    $\Id \in \Hom_\g(V, V)$). In contrast, $V \otimes V \cong S^2V \oplus \Lambda^2V$
    always.
\end{enumerate}
\end{exercise}

\begin{exercise}
Let $\g = \so_n$ with $n = 3$ or $n \geq 5$, or $\g = \spl_{2l}$ with $2l = n$
and $l \geq 2$, and let $V = \C^n$.
\begin{enumerate}
  \item Compute the highest weight of this representation.
  \item $V \cong V^*$ via the form defining $\g$, so $V \otimes V$ has at least
    three components (since it must contain the trivial representation). Show
    that it has exactly three, describe them, and find the highest weights.
\end{enumerate}
\end{exercise}

We will now devote ourselves to proving Theorem~8.3, classifying all the
representations of semisimple Lie algebras.

Let $\g$ be any Lie algebra with a non-degenerate invariant bilinear form
$(\,,\,)$, e.g., $\g$ semisimple and $(\,,\,) = (\,,\,)_{\ad}$. Let
$x_1, \dots, x_N$ be a basis of $\g$, and $x^1, \dots, x^N$ the dual basis (so
$(x_i, x^j) = \delta_{ij}$).

\begin{definition}
$\Omega = \sum x_ix^i$ is the \emph{Casimir} of $\g$.
\end{definition}

\begin{remark}
For the moment, $\Omega$ is an operator on any representation of $\g$.
However, we will in a moment define the ``universal enveloping algebra of
$\g$'', and that is where $\Omega$ naturally lives.
\end{remark}

\begin{lemma}
For all $x \in \g$, we have $[\Omega, x] = 0$.
\end{lemma}

\begin{proof}[Proof 1, coordinate-based]
$[\Omega, x] = [\sum x_ix^i, x]$. Since $\ad x$ is a derivation, we rewrite
this as $\sum x_i[x^i, x] + \sum [x_i, x]x^i$. Write
$[x^i, x] = \sum_j a_{ij}x^j$ and $[x_i, x] = \sum_j b_{ij}x_j$. Then
\begin{align*}
  a_{ij} &= ([x^i, x], x_j) = ([x_j, x^i], x),\\
  b_{ij} &= ([x_i, x], x^j) = ([x^j, x_i], x) = -a_{ji}.
\end{align*}
Thus,
\[
  [\Omega, x] = \sum_{i,j} a_{ij}x_ix^j + \sum_{i,j} b_{ij}x_jx^i = 0. \qedhere
\]
\end{proof}

\begin{proof}[Proof 2, without coordinates]
We have maps of $\g$-modules
$\C \hookrightarrow \End(\g) \cong \g \otimes \g^* \cong \g \otimes \g$ (the
first map is via $\lambda \mapsto \lambda\Id$, i.e.\
$1 \mapsto \sum x_i \otimes x^i$; the last isomorphism is via the bilinear
form). Since $\g$ acts on $V$, we have a map of $\g$-modules
$\g \to \End(V)$. Finally, we have a map of $\g$-modules
$\End(V) \otimes \End(V) \to \End(V)$ by multiplication. This finally gives
\[
  \C \hookrightarrow \g \otimes \g^* \cong \g \otimes \g \to \End(V) \otimes \End(V) \to \End(V)
\]
sending $1 \mapsto \Omega$. Therefore, $\Omega$ generates a trivial
$\g$-submodule of $\End(V)$, which is equivalent to saying $[\Omega, x] = 0$
for all $x \in \g$.
\end{proof}

Now let $\g$ be semisimple, $\g = \at \oplus \bigoplus_{\alpha \in R} \g_\alpha$,
and $(\,,\,)_{\ad}$ the Killing form. Choose a basis $u_1, \dots, u_l$ of
$\at$, $x_\alpha$ of $\g_\alpha$; choose the dual basis $u^1, \dots, u^l$ of
$\at$ and $x^{-\alpha}$ of $\g_{-\alpha}$ (so that $(x_\alpha, x^{-\alpha}) = 1$).
Normalise $x_\alpha$ so that $(x_\alpha, x_{-\alpha})_{\ad} = 1$, so that
$x^{-\alpha} = x_{-\alpha}$. Then $[x_\alpha, x_{-\alpha}] = \nu^{-1}(\alpha)$
by the structure theorem. Now,
\[
  \Omega = \sum u_iu^i + \sum_{\alpha \in R^+} (x_\alpha x_{-\alpha} + x_{-\alpha}x_\alpha)
         = \sum u_iu^i + 2\sum_{\alpha \in R^+} x_{-\alpha}x_\alpha + \sum_{\alpha \in R^+} \nu^{-1}(\alpha).
\]

\begin{definition}
$\rho = \frac12 \sum_{\alpha \in R^+} \alpha$.
\end{definition}

\begin{remark}
Of course, there will be an exercise to compute $\rho$ for all the simple Lie
algebras. However, we will state this exercise later, when it's more obvious
what $\rho$ is.
\end{remark}

\begin{remark}
If you like algebraic geometry, this is the square root of the anticanonical
bundle on the flag variety.
\end{remark}

Then
\[
  \Omega = \sum u_iu^i + 2\nu^{-1}(\rho) + 2\sum_{\alpha \in R^+} x_{-\alpha}x_\alpha.
\]

\begin{lemma}
Let $V$ be a $\g$-module, $v \in V$ a singular vector with weight $\lambda$
(i.e.\ $\fn^+v = 0$ and $tv = \lambda(t)v$). Then
$\Omega v = \bigl(|\lambda + \rho|^2 - |\rho|^2\bigr)v$.
\end{lemma}

\begin{proof}
Apply the most recent version of $\Omega$:
\[
  \Omega v = \Bigl(\sum \lambda(u_i)\lambda(u^i) + \lambda\bigl(2\nu^{-1}(\rho)\bigr)\Bigr)v
    + 2\sum_{\alpha \in R^+} x_{-\alpha}x_\alpha v
  = \bigl((\lambda, \lambda) + 2(\lambda, \rho)\bigr)v. \qedhere
\]
\end{proof}

In particular, if $V$ is irreducible and has a singular vector of weight
$\lambda$, then $\Omega$ acts on $V$ by $(\lambda, \lambda) + 2(\lambda, \rho)$,
by Schur's lemma.
